% id: 142-11
% book: 베이직쎈 미적분1 · 09 정적분의 활용
% section: 유형 090 곡선과 접선으로 둘러싸인 도형의 넓이
% figure: tikz:fig-142-11
% answer: 
% answer_source: 
% variant_level: 0 (원본 전사)
% 이 파일은 build.mjs 가 items.json 에서 만든다. 고칠 때는 items.json 을 고친다.
\smash{\raisebox{1.5mm}{\dmkichul{베이직쎈 미적분1 142쪽 11번}}}
다음은 닫힌구간 $[0{,}\mkern3mu\mathopen{}1]$에서 곡선 $y=x^2+1$과 이 곡선 위의 점 $(1{,}\mkern3mu\mathopen{}2)$에서의 접선 및 $y$축으로 둘러싸인 도형의 넓이를 구하는 과정이다. ㈎$\sim$㈒에 알맞은 것을 구하시오. \exprbox{\begin{minipage}{0.96\linewidth}\raggedright $y=x^2+1$에서 \quad $y'=\text{\fbox{\footnotesize ㈎}}$\par\smallskip 곡선 위의 점 $(1{,}\mkern3mu\mathopen{}2)$에서의 접선의 기울기는 \fbox{\footnotesize ㈏} 이므로 접선의 방정식은\par\smallskip \hspace*{1.5em}$y-2=\text{\fbox{\footnotesize ㈏}}(x-1)$ \qquad $\therefore y=\text{\fbox{\footnotesize ㈐}}$\par\smallskip 따라서 구하는 넓이는\par\smallskip \hspace*{1.5em}$\displaystyle\int_{0}^{1} (\text{\fbox{\footnotesize ㈑}}-2x)\mkern3mu dx$\par\smallskip \hspace*{1.5em}$\displaystyle =\Big[\dfrac{1}{3}x^3+x-x^2\Big]_{0}^{1}$\par\smallskip \hspace*{1.5em}$=\text{\fbox{\footnotesize ㈒}}$\end{minipage}}
\par\smallskip\begin{center}\resizebox{\ifdim\width>\linewidth\linewidth\else\width\fi}{!}{% 142-11(142쪽) — 곡선 $y=x^2+1$ 과 점 $(1,\,2)$ 에서의 접선 $y=2x$ 및 $y$축으로 둘러싸인 색칠한 부분. 스캔 화질이 낮아 TikZ 로 다시 그림.
% 원문 풀이 상자 안의 그림 그대로 옮긴다: 포물선 · 접선 · 색칠 · 점 (1,2) 로 가는 안내 점선 둘 · 라벨 y=x^2+1, y=2x, 1, 2, O, 1, x, y. 원문에 없는 값(넓이)은 적지 않는다.
\begin{tikzpicture}[x=0.46cm, y=0.46cm, line width=0.6pt]
  \fill[gray!25] plot[domain=0:1, samples=60, smooth] (\x, {(\x)*(\x)+1}) -- plot[domain=1:0, samples=60, smooth] (\x, {2*(\x)}) -- cycle;
  \draw[->, >=stealth, line width=0.7pt] (-2.0,0) -- (2.4,0) node[below=1pt, font=\small] {$x$};
  \draw[->, >=stealth, line width=0.7pt] (0,-0.95) -- (0,4.55) node[left=1pt, font=\small] {$y$};
  \draw[densely dashed, line width=0.4pt] (0,2) -- (1,2);
  \draw[densely dotted, line width=0.4pt] (1,2) -- (1,0);
  \draw[domain=-1.55:1.5, samples=100, smooth] plot (\x, {(\x)*(\x)+1});
  \draw (-0.45,-0.9) -- (1.75,3.5);
  \node[anchor=east, inner sep=1pt, font=\small] at (-0.75,3.85) {$y=x^2+1$};
  \node[anchor=west, inner sep=1pt, font=\small] at (1.15,3.85) {$y=2x$};
  \node[left=2pt, font=\small] at (0,2) {$2$};
  \node[left=2pt, font=\small] at (0,1) {$1$};
  \node[below right=-2pt and -2pt, font=\small] at (0,0) {$\pt{O}$};
  \node[below=1pt, font=\small] at (1,0) {$1$};
\end{tikzpicture}
}\end{center}
